% !TEX root = ../notes_template.tex
\chapter{Clinical Skills Checklist}\label{chp:skills_checklist}

%\addcontentsline{toc}{chapter}{Week 10: Clinical Skills Checklist}

This checklist summarizes the key psychomotor skills and reasoning tasks practiced throughout the Neuromuscular \& Vascular Lab. It is designed to help you track your progress, reinforce consistency, and prepare for your practical exam.

You may use this checklist during open labs, peer practice, or recitation sessions. All items must eventually be demonstrated to an instructor or lab assistant for formal sign-off. Each item has two checkboxes: one for peer or self-confirmation, and one for official instructor verification.

\subsection*{Week 1: Movement, MMT, and MLT Foundations}

\begin{itemize}
  \item [$\square$] Performed and interpreted SFMA Top Tier patterns, with identification of major muscle contributors.
  \item [$\square$] Conducted gravity-resisted MMT for hip flexion and shoulder abduction, with break test and HHD measurement.
  \item [$\square$] Performed gravity-eliminated MMT for one assigned motion, using appropriate positioning and stabilization.
  \item [$\square$] Documented primary movers, observed strength, and compensation patterns for MMT scenarios.
  \item [$\square$] Conducted modified Thomas test and interpreted rectus femoris length, including R1, R2, and observed compensations.
  \item [$\square$] Conducted straight leg raise (SLR) for hamstrings, with R1, R2, and interpretation of sequencing and compensations.
\end{itemize}

\subsection*{Week 2: Reflexes, Pulses, and Rehabilitative Imaging}

\begin{itemize}
  \item [$\square$] Performed and graded deep tendon reflexes (DTRs) for biceps and patellar tendons, using the 0–4+ scale.
  \item [$\square$] Interpreted DTR findings in relation to muscle weakness and motor neuron involvement.
  \item [$\square$] Identified musculoskeletal tissue types (muscle, tendon, nerve, bone, vessel) on ultrasound using proper LAX and SAX orientation.
  \item [$\square$] Applied the “5 Ps” of scanning and demonstrated effective use of probe motion techniques (heel-toe, fanning, sliding, rotation).
  \item [$\square$] Acquired and interpreted B-mode ultrasound images of the patellar and biceps tendons in both LAX and SAX views.
  \item [$\square$] Visualized and confirmed the radial artery using ultrasound, assessing pulsatility and compressibility.
  \item [$\square$] Compared palpated radial pulse with ultrasound visualization, noting symmetry and vessel location.
\end{itemize}


\subsection*{Week 3: Shoulder Complex}
\begin{itemize}
  \item [$\square$] Performed gravity-resisted and gravity-eliminated MMT for glenohumeral and scapular muscles.
  \item [$\square$] Identified the innervation and nerve root(s) of all muscles tested during MMT.
  \item [$\square$] Used HHD to record force output for one GH break test.
  \item [$\square$] Conducted and interpreted MLT for pectoralis major, pectoralis minor, latissimus dorsi, and teres major, including R1, R2, and compensations.
  \item [$\square$] Acquired static ultrasound images of the supraspinatus tendon in SAX and LAX.
  \item [$\square$] Performed dynamic ultrasound to assess AC joint space during shoulder abduction.
  \item [$\square$] Palpated and imaged the axillary pulse, confirming vascular characteristics via ultrasound.
\end{itemize}

\subsection*{Week 4: Elbow, Forearm, Wrist \& Hand}
\begin{itemize}
  \item [$\square$] Performed gravity-resisted and gravity-eliminated MMT for elbow flexion, extension, supination, and pronation with correct stabilization and substitution awareness.
  \item [$\square$] Used hand-held dynamometry (HHD) to quantify force output for either elbow flexion or extension.
  \item [$\square$] Performed and graded deep tendon reflexes for the biceps (C5–C6) and triceps (C7) tendons.
  \item [$\square$] Conducted MMT for extrinsic hand function (grip, finger extension, FDP) to screen radial and ulnar nerve involvement.
  \item [$\square$] Compared grip strength with wrist in flexed vs. extended positions to evaluate influence of active and passive insufficiency.
  \item [$\square$] Conducted MMT for intrinsic hand muscles (APB, adductor pollicis, 1st dorsal interosseous, lumbricals) for differential median vs. ulnar nerve screening.
  \item [$\square$] Identified and interpreted upper extremity nerve entrapment sites from elbow to wrist using anatomical landmarks.
  \item [$\square$] Acquired and interpreted ultrasound images of the median and ulnar nerves in both SAX and LAX views using dynamic validation techniques.
  \item [$\square$] Palpated and imaged the brachial and radial arteries using ultrasound to confirm pulsatility and compressibility.
\end{itemize}

\subsection*{Week 5: Pelvis \& Hip}
\begin{itemize}
  \item [$\square$] Interpreted SFMA Top Tier patterns (MSF, MSE, MSR, SLS, ADS) with attention to pelvic-hip coordination and regional interdependence.
  \item [$\square$] Performed gravity-resisted and gravity-eliminated MMT for major hip movements: flexion, extension, abduction, adduction, internal rotation, external rotation.
  \item [$\square$] Identified peripheral nerve and spinal root(s) for each primary muscle tested in MMT.
  \item [$\square$] Repeated the Straight Leg Raise (SLR) to assess hamstring extensibility and compare to Week 1 performance.
  \item [$\square$] Performed the Modified Thomas Test to assess iliopsoas and rectus femoris length using R1/R2 and compensatory pattern recognition.
  \item [$\square$] Conducted the Ober test for TFL/IT band extensibility with appropriate stabilization.
  \item [$\square$] Visualized the femoral neurovascular bundle using ultrasound and confirmed anatomical landmarks.
\end{itemize}


\subsection*{Week 6: Knee}

\begin{itemize}
  \item [$\square$] Interpreted SFMA patterns (MSF, MSE, MSR, SLS, ADS) with attention to tibiofemoral mechanics and regional interdependence.
  \item [$\square$] Performed MMT for knee flexion and extension, identifying key muscle groups and neural innervation.
  \item [$\square$] Used hand-held dynamometry to compare quadriceps and hamstrings strength and analyze functional ratio.
  \item [$\square$] Performed the Nordic Hamstring Exercise as a screen for eccentric knee flexor control.
  \item [$\square$] Performed the patellar tendon reflex and interpreted findings in relation to L3–L4 and femoral nerve function.
  \item [$\square$] Conducted the Modified Thomas Test to evaluate rectus femoris length with R1 and R2 documentation.
  \item [$\square$] Acquired RUSI images of the distal quadriceps, patellar tendon, and suprapatellar pouch for effusion assessment.
  \item [$\square$] Palpated and imaged the popliteal artery, confirming pulsatility and compressibility with ultrasound.
\end{itemize}

\subsection*{Week 7: Ankle \& Foot}

\begin{itemize}
  \item [$\square$] Interpreted SFMA patterns (ADS, SLS, MSF, MSR) with attention to dorsiflexion, pronation, and medial arch mechanics.
  \item [$\square$] Performed MMT for ankle and foot actions: dorsiflexion, plantarflexion, inversion, eversion, toe flexion.
  \item [$\square$] Identified the peripheral nerve and root level for each tested muscle group.
  \item [$\square$] Performed the Achilles tendon reflex and interpreted findings in relation to S1 and tibial nerve integrity.
  \item [$\square$] Used RUSI to image the Achilles tendon in LAX and SAX at three locations: insertion, mid-substance, myotendinous junction.
  \item [$\square$] Performed introductory dynamic scanning of the lateral ankle to visualize ATFL and CFL during inversion/plantarflexion.
  \item [$\square$] Palpated and imaged the dorsalis pedis or posterior tibial pulse, confirming vessel identity with ultrasound.
\end{itemize}

\subsection*{Week 8: Lumbopelvic}
\begin{itemize}
  \item [$\square$] Applied SFMA Top Tier movement logic to identify lumbopelvic control deficits.
  \item [$\square$] Performed trunk flexion and extension MMT using standardized criteria.
  \item [$\square$] Executed leg-lowering test to assess trunk motor control with angle of loss recorded.
  \item [$\square$] Identified primary trunk muscles with associated nerve and root innervation.
  \item [$\square$] Acquired and interpreted ultrasound images of TrA and multifidus.
  \item [$\square$] Analyzed muscle recruitment quality, symmetry, and substitution patterns.
\end{itemize}

\subsection*{Week 9: Cervical \& Quarter Screens}

\begin{itemize}
  \item [$\square$] Performed MMT for cervical flexion, extension, rotation, and sidebending.
  \item [$\square$] Identified innervation and root levels for key cervical musculature.
  \item [$\square$] Conducted the Upper Quarter Screen: myotomes, dermatomes, and reflexes.
  \item [$\square$] Conducted the Lower Quarter Screen: myotomes, dermatomes, and reflexes.
  \item [$\square$] Interpreted screen findings in the context of neuroanatomy and regional interdependence.
\end{itemize}